\section{Introduction}

Models for high stakes decision making have ethical and legal needs to demonstrate a lack of discrimination with respect to protected classes \cite{Munoz2016,Barocas:2016aa}.
Examples of such decisions include employment and compensation \cite{Conway1983,Greene1984}, university admissions \cite{Bickel1975}, and sentence and bail setting \cite{Berk2017,Chouldechova2017,Dressel2018}.
Another example relevant to the financial services industry is credit decisioning \cite{Chen2018}, which is a classification problem where these ethical concerns are enshrined in concrete regulatory compliance requirements.
Credit decisions must be shown to comply with a myriad of federal and state fair lending laws, some of which are summarized in \cite{Chen2018}%
\footnote{In this paper, we restrict our discussion and citation of applicable laws to those of the United States of America.}.
Some of these laws define protected classes, such as race and gender,
where discrimination on the basis of a customer's membership in these classes is prohibited. Table \ref{table:classes} summarizes the protected classes defined by the Fair Housing Act (FHA) \cite{fha} and Equal Credit Opportunity Act (ECOA) \cite{ecoa}. 

%compliance need creates a conundrum In this paper, we analyze the biases inherent in using the Bayesian Improved Surname Geocoding (BISG) model\cite{Kevin:2006aa} for imputing race labels, which is known to be used by regulators \cite{zhang2016assessing,CFPB:proxy2014}, and provide a statistical analysis of how imputation bias interacts with downstream measures of disparity.

%\subsection{The regulatory needs for fairness in credit decisioning}

%such as the
%Equal Credit Opportunity Act (ECOA) \cite{ecoa},
%Fair Credit Reporting Act (FCRA) \cite{fcra},
%Fair and Accurate Credit Transactions Act (FACTA) \cite{facta},
%Fair Housing Act (FHA) \cite{fha},
%Fair and Equal Housing Act (FEHA) \cite{feha},
%and
%Consumer Credit Protection Act (CCPA) \cite{ccpa}.
%(see Table~\ref{table:classes})

\begin{table}[h]

\caption{\label{table:classes} Protected classes defined under fair lending laws.}

\begin{tabular}{l c c}
\toprule
Law & FHA\cite{fha}\! & ECOA\cite{ecoa}\! \\
\midrule
% \hline\hline
age & & X \\
color & X & X \\
disability & X & \\
exercised rights under CCPA & & X \\
familial status (household composition)\!\!& X & \\
gender identity & X & \\
marital status (single or married)& & X \\
national origin & X & X \\
race  & X & X \\
recipient of public assistance & & X \\
religion & X & X \\
sex & X & X \\
% \hline
\bottomrule
\end{tabular}

\end{table}


%\paragraph{Challenges in evaluating unfairness in credit decisioning}
%In practice, evaluating the fairness of decisioning models is challenging for both legal and statistical reasons.
%For regulated decisions, there are two major legal theories of discrimination:
%
%\begin{description}
%
%\item[disparate treatment]\hspace{-1.3mm} \cite{Zimmer:1996aa}, informally, intentionally treating an individual differently on the basis of membership in a protected class, and
%
%\item[disparate impact]\hspace{-1.3mm} \cite{Rutherglen:1987aa}, informally, adversely affecting members of one protected class more than another even if by an ostensibly neutral policy.
%% discriminating against any protected class as a resulting from implementing of a facially neutral policy.
%\end{description}
%%
%Different theories of discrimination may apply for different legal situations.
%Disparate impact as a theory of discrimination under FHA has been confirmed by the Supreme Court; however, there is some uncertainty as to whether disparate impact is a proper theory of discrimination under ECOA \cite{Cubita:2006aa,Avery:2012aa,Ropiequet:2014aa,McDonald2015:aa}.
%While several courts and regulators believe that disparate impact should be considered under ECOA \cite{OCC:2017,CFPB:2018}, the Bureau of Consumer Financial Protection (CFPB) has signaled that it may issue guidance in the near future that will limit the use of disparate impact claims under ECOA \cite{CFPB:2018-05}.
%
%Furthermore, the nature of statistical evidence considered in compliance assessments can vary contextually.
%For example, the ``four-fifths rule'' is used in assessing adverse impact with regard to racial discrimination in US employment law \cite{eeoc:1978}; however, no such rules of thumb exist for assessing adverse impact in fair lending \cite{OCC:2017,regulationb,CFPB:2018}.
%As a result, the admissibility of any statistical test of fairness to assess compliance must be considered in the appropriate legal, ethical and regulatory context.
%This context sensitivity complicates the practical choices of metrics of fairness \cite{Chouldechova2017,Kusner2017,Zliobaite2017}, on top of the known impossibility of satisfying mutually contradictory measures simultaneously \cite{Conway1983,Greene1984,Friedler2016,Kleinberg2017}.

\subsection{Assessing fairness with unknown protected class membership}

When demonstrating that credit decisions comply with these fair lending laws,
we sometimes run into situations where fairness and bias assessments must be done on populations without knowing their memberships in protected classes, because it is illegal or operationally difficult to do so.
For example, credit card and auto loan companies must demonstrate that the way they extend credit is not racially discriminatory, yet are not allowed to ask applicants what race they are when they apply for credit.\footnote{Lenders may ask applicants to self-identify in a voluntary basis, with the understanding that the answer will not affect the outcome of the application and that the information is collected for compliance assessment only \cite[12 CFR \S 1002.5(b)]{regulationb}.}
Similarly, health plans can only solicit race and ethnicity information for new members but cannot obtain the same information for existing members \cite{elliott2008new}.
% In these cases, even regulators do not have access to the protected attributes information, yet they still need to evaluate the outcome disparity with respect to these protected attribtues.
Given the lack of secure protocols that permit disparity evaluation with encrypted protected classes \cite{kilbertus2018blind}, disparate impact assessments for these situations have to impute the mostly (or entirely) missing labels corresponding to the protected class, usually by relying on observed proxy variables that can predict class memberships.
 The imputed protected classes are then used by regulators in assessing disparate impact (but they are not allowed to be used in decision making).
 Generally, any model that imputes the missing protected attribute value based on other, observed variables is known as a \emph{proxy model}, and such a model that is based on predicting conditional class membership probabilities is known as a \emph{probabilistic proxy model}.
% the probabilities for the value of missing protected attributes.

%\paragraph{Proxy models used in practice}

For example, for assessing adverse action with regard to race in credit decisions, regulators like the Bureau of Consumer Financial Protection (BCFP)\footnote{Formerly the Consumer Financial Protection Bureau (CFPB). Citations and references reflect the name at time of publication.}  have been known in the past to use a probabilistic proxy model to impute the customers' unknown race labels \cite{CFPB:proxy2014}.
They used a na\"ive Bayes classifier, the Bayesian Improved Surname Geocoding (BISG) method, to predict the probability of race membership given the customer's surname and address of residence \cite{Elliott:2009aa}.
Specifically, assuming that surname and location are statistically independent given race, BISG uses Bayes's rule to
compute race membership probabilities from
the conditional distributions of surname given race and of location given race as inferred by census data.
This methodology \cite{CFPB:proxy2014} notably supported a \$98 million fine against a major auto loan lender \cite{CFPB:ally2013}.
This case generated some controversy \cite{Koren:2016aa,CFS:2015,CFS:2016,CFS:2017,CFPB:2018-05}, in part due to empirical findings that the amount of disparate impact estimated by BISG appears to overestimate true disparities \cite{Baines:2014aa,zhang2016assessing}.
However, the cause for this overestimation phenomenon is unknown, as is whether overestimation is to be expected always, or whether or not underestimation of disparate impact is also possible.
This observation forms the motivation for our current work, which is broadly applicable to any fairness assessment where an unobserved protected class must be imputed using a proxy model.
The aim is not to criticize the use of proxy models in general, but rather to provide a more informed analysis of the statistical biases inherent in any assessment where membership in protected classes must be imputed.


\paragraph{Main results}

This paper investigates the bias in estimating demographic disparity (Definition \ref{demographic-disparity}) when a proxy model is employed to impute a protected class.
% We call any such imputation model a \textit{probabilistic proxy}.
We present the first theoretical results describing when the use of proxy models can lead to biased estimates of outcome disparity,
which explains the overestimates observed in the past, and also offers insights on the practical use of proxy models.
More specifically, our key contributions are:

\begin{enumerate}
\item
We derive the (statistical) bias for the commonly used thresholded estimator, where a label is assigned only if the proxy model predicts a label with probability exceeding a predefined threshold \cite{CFPB:proxy2014,Baines:2014aa,zhang2016assessing} (Definition \ref{def:thresholded_estimator}, Theorem \ref{thm: demographic-est-race-bias}). We decompose its bias into multiple sources, which gives a set of interpretable conditions under which the thresholded estimator can over- or underestimate the outcome disparity.

\item
We present a new weighted estimator for demographic disparity (Definition \ref{def:weighted_estimator}) that uses soft classification based on proxy model outputs as opposed to hard imputation. We derive its bias (Theorem \ref{thm: weighted}) and find that the weighted estimator has only one bias source.
% We show evidence that it tends to underestimate demographic disparity, and thus can be used with the thresholded estimator to bracket the true demographic disparity (Section \ref{sec:numerics}).

% \item
% We provide a set of interpretable conditions under which these two estimators can overestimate or underestimate the true demographic disparity, and categorize the different sources of bias.

\item
We validate our results on a public mortgage data set, using geolocation as the sole variable in a proxy model for race.
We identify the specific source of bias that can account for the overestimation of the thresholded estimator, which can explain the overestimation bias of using proxy methods observed in previous literature.
%The result extends to any proxy model that estimate the protected attribute based on predictors correlated with the decision outcome (e.g., the geolocation used in BISG is known to be correlated with socioeconomic status that can affect lending outcome).
\item
We discover that the estimation bias is sensitive to the threshold used in class imputation based on the proxy model. This shows the intrinsic limitation of the thresholded estimator.
\end{enumerate}



% \paragraph{Literature on probabilistic proxies}

% Fiscella and Fremont conducted a comprehensive review on the race proxy methods for health care disparity assessment, finding that surname analysis and geocoding analysis have different predictive powers for different races, and they suggest a non-Bayesian hybrid approach that combines both surname and geocoding information for race imputation \cite{fiscella2006use}. formaliedcame up with the BISG method which uses Bayes rule to combine geolocation and surname information and showed that BISG is more accurate in race classification than previous race proxy methods on a health plan dataset. CFPB \cite{CFPB:proxy2014} also validates the same conclusion on a mortgage dataset and introduces BISG in fair lending assessment. Voicu \cite{voicu2018using} augmented the BISG method with first name information, which considerably improves the race classification accuracy for African American for whom previous methods have high classification error rate. Zhang \cite{zhang2016assessing}, Baines and Courchane \cite{Baines:2014aa} considered using the BISG proxy method in evaluating mortgage loan disparity and found that using BISG method leads to overestimates of the outcome disparity.

%\textbf{Organization}. In Section 2, we formulate the problem of using proxy for unknown protected attribute when evaluating the outcome disparity, introduce the demographic disparity measure, the thresholded estimator and the weighted estimator. We also present two toy lending examples to develop the intuition on the estimation bias of the thresholded estimator. In Section 3, we theoretically analyze the bias of these two estimators and discuss scenarios under which they overestimate or underestimate the disparity measure. In section 4, we back up our theoretical result using both numerical simulation and a mortgage dataset. We finally conclude with a discussion about the limitations of using proxy methods.

% \subsection{Toy example: pens in drawers}

% Suppose you clean up a bunch of loose blue and pens around a conference room. At a glance you don't know what color ink the pen has, so you test out if they still write and sort them into drawers in a cabinet. If they write well you put them in drawer \#1, and if they don't write well you throw them in drawer \#2 (or in the trash). When someone needs a pen, they can grab it from drawer \#1 and know that they will have a pen that works well.

% Suppose the number of blue and red pens in each drawer is captured in Table \ref{table:pens1}. We can calculate if there's a discrepancy, $D$, between working pens for each ink color by computing:
% \begin{equation*}
% \begin{aligned}
% D &= \Big(\parbox{9em}{average of red pens that still work}\Big) - \Big(\parbox{9em}{average of blue pens that still work}\Big) \\
%   &= \frac{8 \text{working red}}{12 \text{total red}} - \frac{2 \text{working blue}}{8 \text{total blue}} \\
%   &= 0.417
% \end{aligned}
% \end{equation*}

% This discrepancy appears because there is a positive correlation between red ink color and "writes well" and a negative correlation between blue ink and "writes well".


% Commenting out this version of the toy example
% \subsection{Toy example: pens in drawers}
%
% Suppose you clean up a bunch of loose blue and pens around a conference room. At a glance you don't know what color ink the pen has, so you test out if they still write and sort them into drawers in a cabinet. If they write well you put them in drawer \#1, and if they don't write well you throw them in drawer \#2 (or in the trash). When someone needs a pen, they can grab it from drawer \#1 and know that they will have a pen that works well.
%
% Suppose the number of blue and red pens in each drawer is captured in Table \ref{table:pens1}. We can calculate if there's a discrepancy, $D$, between working pens for each ink color by computing:
% \begin{equation*}
% \begin{aligned}
% D &= \Big(\parbox{9em}{average of red pens that still work}\Big) - \Big(\parbox{9em}{average of blue pens that still work}\Big) \\
%   &= \frac{8 \text{working red}}{12 \text{total red}} - \frac{2 \text{working blue}}{8 \text{total blue}} \\
%   &= 0.417
% \end{aligned}
% \end{equation*}
%
% This discrepancy appears because there is a positive correlation between red ink color and "writes well" and a negative correlation between blue ink and "writes well".
%
%
% \begin{table}
%   \caption{Working pens sorted into different drawers}
%   \label{table:pens1}
%   \begin{tabular}{c | c | c}
%     \hline\hline
%     & \multicolumn{2}{c}{True pen color}\\
%     \hline
%     Drawer & red & blue \\
%     \hline
%     \#1 & 8 & 2 \\
%     \#2 & 4 & 6 \\
%     \hline
%   \end{tabular}
% \end{table}
%
%
% \begin{table}
%   \caption{Guessing pen color based on drawer}
%   \label{table:pens2}
%   \begin{tabular}{c | c | c | c }
%     \hline\hline
%     & \multicolumn{3}{c}{Guessed pen color}\\
%     \hline
%     Drawer & red & blue & ? \\
%     \hline
%     \#1 & 10 & 0 & 0 \\
%     \#2 & 0 & 0 & 10\\
%     \hline
%   \end{tabular}
% \end{table}
%
% If we can't check each individual pen's color, we can estimate the discrepancy by guessing at pen color for an entire drawer. For example, let's say that if more than 70\% of the pens in a drawer are a certain color, we'll treat all the pens in that drawer as a single color, and if not we won't specify the color and call it a "?". Applying this threshold to our counts in Table \ref{table:pens1}, we get the guess numbers in Table \ref{table:pens2} and the discrepancy $D_\text{estimated}$:
%
% \begin{equation*}
% \begin{aligned}
% D_{\text{estimated}} &= \Big(\parbox{9em}{average guessed-red pens that still work}\Big) - \Big(\parbox{9em}{average guessed-blue pens that still work}\Big) \\
%   &= 1 - 0 \\
%   &= 1.0
% \end{aligned}
% \end{equation*}
%
% Using the drawer number to guess pen colors produces a discrepancy overestimate compared to the discrepancy produced by using the true pen colors. The difference between these two measures is a bias introduced by use of a probablistic proxy for an unknown or unmeasured attribute. By applying the same steps as this toy example, we can analyze disparity measures used when discussing model fairness. In the following sections, we will identify correlation conditions between outcomes and the probablistic proxy to explain the behavior of the bias and indicate when it creates an disparity overestimate for a given dataset.
% End commenting out this version of the toy example

%
%
% \subsection{Toy example}
% We explore a scenario to demonstrate that a disparity measure estimate can be different from the true value when using a proxy for an unknown protected attribute. We simplify the scenario by considering only two neighborhoods, with different average-income levels, and two groups of people, left handed and right handed, in Table \ref{table:toy-true}.

% \begin{table}
%   \caption{True handedness by neighborhood}
%   \label{table:toy-true}
%   \begin{tabular}{c | c | c}
%     \hline\hline
%     & \multicolumn{2}{c}{true handedness}\\
%     \hline
%     income level & right & left \\
%     \hline
%     high & 80 & 20 \\
%     low & 60 & 40 \\
%     \hline
%   \end{tabular}
% \end{table}

% We compute the demographic disparity between the right handers and left handers for a simplified credit offer based on income:

% \begin{equation*}
% \begin{aligned}
% D_\text{true} &= \Big(\parbox{9em}{\% right-handers with credit offer}\Big) - \Big(\parbox{9em}{\% left-handers with credit offer}\Big) \\
%   &= \frac{80 \text{offers}}{80+60 \text{total}} - \frac{20 \text{offers}}{20+40 \text{total}} \\
%   &= 0.238
% \end{aligned}
% \end{equation*}

% The disparity appears in this example because, based on the group populations in Table \ref{table:toy-true}, there is a positive correlation between being right-handed and living in the higher-income neighborhood.

% \paragraph{Introducing a threshold-based probablistic proxy for right- and left-handers}
% Under circumstances where it is cost-prohibitive, illegal, or unethical to collect information about a person's handedness, we may still want to estimate the disparate impact of the credit offer. We may be able to sample, or use data aggregates, to find out the proportion of right-handers in a particular neighborhood, even though we don't know that information at the individual level. One proxy we can apply is to say that when a neighborhood is above some threshold, say, 70\%, then we'll assign everyone in that neighborhood to the majority group, and when below the threshold we'll assign that neighborhood as "unknown". Applying this proxy to the data in Table \ref{table:toy-true} yields Table \ref{table:toy-estimate}.

% \begin{table}
%   \caption{Proxy handedness by neighborhood (70\% threshold)}
%   \label{table:toy-estimate}
%   \begin{tabular}{c | c | c | c }
%     \hline\hline
%     & \multicolumn{3}{c}{proxy handedness}\\
%     \hline
%     income level & right & left & unknown \\
%     \hline
%     high & 100 & 0 & 0 \\
%     low & 0 & 0 & 100 \\
%     \hline
%   \end{tabular}
% \end{table}

% We can carry out the disparity calculation to get an estimate based on the proxy values:

% \begin{equation*}
% \begin{aligned}
% D_\text{estimate} &= \Big(\parbox{9em}{\% est. right-handers with credit offer}\Big) - \Big(\parbox{9em}{\% est. left-handers with credit offer}\Big) \\
%   &= \frac{100 \text{offers}}{100 \text{total}} - 0 \\
%   &= 1.0
% \end{aligned}
% \end{equation*}

% Using the proxy for handedness produces a disparity overestimate compared to the true disparity, i.e. $D_\text{true} < D_\text{estimate}$. The difference between these two measures is a bias introduced by use of a probablistic proxy for an unknown or unmeasured attribute. By applying the same steps as this toy example, we can analyze disparity measures used to discuss model fairness. In the following sections, we will identify correlation conditions between outcomes and the probablistic proxy to explain the behavior of the bias and indicate when it creates an disparity overestimate for a given dataset.

%
% \begin{itemize}
%   \item{$\hat{Y} = 1$ is an offer of credit}
%   \item{$Z_1$ and $Z_2$ are geographic neighborhoods with above- and below-average socioeconomic status, respectively}
%   \item{$A$ is the true race of an individual}
%   \item {$\hat{A}$ is the estimated race of an individual}
%   \item{$a$ and $b$ are the race values for a pair of racial groups in a disparate impact analysis, for example, representing a White-Black comparison, or White-Hispanic comparison, etc.}
%   \item{$E(\hat{Y} \mid A = a) - E(\hat{Y} \mid A = b)$ is the demographic disparity}
% \end{itemize}
%
% \paragraph{Scenario 1}
% This scenario assumes the outcome of credit offer is strongly correlated with the socioeconomic status, which is in turn strongly correlated with geographic neighborhood, as summarized in Table \ref{table:scenario1}
%
% \begin{table}
%   \caption{Scenario 1: Outcomes strongly correlated with neighborhood}
%   \label{table:scenario1}
%   \begin{tabular}{c | c | c}
%     \hline\hline
%     Neighbood & $A = a$ & $A = b$ \\
%     \hline
%     $Z_1$ & 30 people (100\%) & 70 people (100\%) \\
%     $Z_2$ & 70 people (0\%) & 30 people (0\%) \\
%     \hline
%   \end{tabular}
% \end{table}
%
% Computing demographic disparity given the ground truth labels gives:
%
% \[ E(\hat{Y} \mid A = a) - E(\hat{Y} \mid A = b) = \frac{30\times1.00}{100} \]

%
% Now we use two simple examples to demonstrate the conditions. In the following examples, we have two locations (rich and poor) and two races (black and white). Each entry in the tables stands for the number of people of a certain race living in a certain location. The numbers in the parentheses represents the credit acceptance rate for people.
%
% Notation:
% - $Y = 1$ gives credit
% - $A = w, b$ for the race
% - $\hat{A} = w, b$ for the estimated race
% - $Z = r, p$ for the locations
% - demographic disparity: $E(Y \mid A = w) - E(Y \mid A = b)$
%
% ## Example 1 for condition (a)(b)
% |      | Black            | White            |
% |------| -----------------|------------------|
% | rich | 30 people (100%) | 70 people (100 %)|
% | poor | 70 people (0%)   | 30 people (0%)   |
%
% - Demographic disparity given the ground truth race lable: $\frac{30*0\% + 70*100\%}{100} - \frac{30*100\% + 70*0\%}{100} = 40\%$
% - According to thresholding rule, we would classify all people in the rich location as white since $P(A = w \mid Z = r) = 0.7 > 0.5$; we classify all people in the poor location as black since $P(A = b \mid Z = p) = 0.7 > 0.5$. So the new table would be:
%
% |      | Estimated Black            | Estimated White            |
% |------| -----------------|------------------|
% | rich | 0 people  | 100 people (100 %)|
% | poor | 100 people (0%)   | 0 people    |
%
% - Demographic disparity based on the estimated race: $E(Y \mid \hat{A} = w) - E(Y \mid \hat{A} = b) = 1 - 0 = 100\%$.
% - This example shows the type I bias captured by (a)(b), i.e., the positive correlation between $Y$ and $P(A = w \mid Z)$ and the negative correlation between $Y$ and $P(A = b \mid Z)$.
%
% ## Example for condition (c)(d)
% |      | Black            | White            |
% |------| -----------------|------------------|
% | rich | 30 people (90%) | 70 people (60 %)|
% | poor | 70 people (20%)   | 30 people (10%)   |
%
% - Demographic disparity given the ground truth race lable: $\frac{70*60\% + 30*10\%}{100} - \frac{30*90\% + 70*20\%}{100} = 4\%$
% - The estimated race is the same as before, but now the acceptance rate is kind of different:
%
% |      | Estimated Black            | Estimated White            |
% |------| -----------------|------------------|
% | rich | 0 people | 30 people (90%), 70 people (60 %)|
% | poor | 70 people (20%), 30 people (10%) | 0 people |
% - Demographic disparity based on the estimated race: $E(Y \mid \hat{A} = w) - E(Y \mid \hat{A} = b) = \frac{30*90\% + 70*60\%}{100} - \frac{30*10\% + 70*20\%}{100} = 52\%$.
% - This example shows the type II bias captured by condition (c)(d), namely within the same $Z$ level (and thus the same probability estimate range), black people have higher average outcome than white people. In practice, I would expect type II bias to be quite small, especially when black people is the minority group so they really need average oucome much higher than their white neighbors to make this type II bias to be noticeable.


% \subsection{Overview}

% We use a proxy for the protected attribute by assigning a value with a threshold cutoff method and a weighted estimator method.

% We compare the bias introduced by the proxy methods for the unknown protected attribute across several fairness measures of disparate impact.

% We provide the first theoretical results explaining when proxy methodologies over- or under-estimate disparity, under a few simple interpretable assumptions. The results generally apply to any proxy methodology when the protected attribute is unknown, not limited to BISG for race. We connect the results to relevant case-studies in education and the financial industry.

% We illustrate our results with simple simulated scenarios and on a public mortgage dataset.

% We discuss the limitations of proxy methods, call for judicious usage of proxy, and advocate collecting the ground truth protected attribute.
